\section{Introduction}
\label{sec:introduction}

Active Directory Certificate Services (ADCS) issues certificates that can be used for domain authentication through PKINIT~\cite{rfc4556}. Different services record issuance and subsequent use. When a certificate template is changed, used and restored between scans, its current configuration no longer describes the conditions under which an earlier certificate was issued. Investigators must therefore connect directory changes, CA records and KDC authentication outcomes.

SpecterOps documented certificate-service escalation techniques~\cite{schroeder2021certified}. Certificate mapping and the installed Windows updates also affect whether a certificate is accepted for an administrative identity~\cite{microsoft2025strongbinding}. We ask what certificate-specific attribution is justified when identifiers or timing evidence are incomplete, late or contradictory.

\subsection{Certificate Issuance and the Evidence Gap}
In an ADCS deployment, certificate issuance is governed by Active Directory objects known as \textit{Certificate Templates} (stored in the LDAP configuration naming context). When an enterprise template exposes an insecure combination of attributes—such as the \texttt{CT\_FLAG\_ENROLLEE\_SUPPLIES\_SUBJECT} bitflag alongside an Extended Key Usage (EKU) permitting Client Authentication (\texttt{1.3.6.1.5.5.7.3.2}) or Smart Card Logon (\texttt{1.3.6.1.5.2.3.4}), without requiring CA administrator approval—any domain principal granted enrollment permissions can request an X.509 certificate specifying an arbitrary Subject Alternative Name (SAN). The attacker can supply the user principal name (UPN) of a Domain Administrator (e.g., \texttt{Administrator@domain.local}), obtain a cryptographically valid certificate signed by the enterprise root CA, and submit that certificate to the Kerberos Key Distribution Center (KDC) via PKINIT to request a Ticket Granting Ticket (TGT), when the KDC accepts the resulting identity mapping. Certificate-binding policy and patch level can prevent this outcome~\cite{microsoft2025strongbinding}.

Configuration assessment and incident reconstruction answer different questions. Certipy reports template and CA risks~\cite{lyak2021certipy}; a positive finding does not by itself show that a certificate was used with elevated privileges. Whole-network provenance systems connect activity across hosts~\cite{liu2026hades,liu2026cpd}, but a certificate-level explanation still needs an issuer-scoped link between issuance and use.

This study examines three issues. First, a graph classifier can learn template attributes that correlate with generated labels while missing the access-control path. Such shortcut behavior is a known learning problem~\cite{geirhos2020shortcut}; our experiments test it under a stated ADCS permission abstraction. Second, restored configuration attributes cannot identify every intervening issuance. Third, the captured CA and KDC records expose different certificate fields. Existing SAN-based CA/KDC queries provide useful operational correlation~\cite{insecurity2023adcs,splunk2026esc1}. We examine what can still be concluded when the required fields are absent or disagree.

\subsection{Research Question and Contributions}
The central question is whether an investigator can make a certificate-specific issuance-to-use claim when configuration has been restored and evidence is incomplete. CertGraph treats identity linkage, temporal overlap and remediation as separate decisions: a compatible identifier does not establish temporal causation, and a risky configuration does not prove administrative certificate use. The contributions are:
\begin{itemize}
\item \textbf{An issuer-scoped evidence contract.} CIPG retains compatible issuance candidates, requires agreement among shared identifiers, distinguishes possible from definite overlap, and revises links when evidence conflicts. A guarded SQL realization tests whether these requirements can be implemented relationally.
\item \textbf{Controlled certificate-use evidence.} Frozen synthetic studies measure candidate retention and definite coverage. Fixed-request interventions and a prospective concurrent workload supply 72 independently outcome-checked certificate cases, including current-patch rejection controls and administrative use after template restoration.
\item \textbf{Complementary remediation analysis.} A sequential, cost-aware interdiction heuristic respects immutable permissions. Evaluation covers 120 graph/budget instances, exhaustive small cases and an equal-query path-preserving projection. Assigned disruption costs and measured cut counts remain distinct from business disruption.
\end{itemize}
An auxiliary graph-learning study tests an optional screening mechanism. Additive and literal bilinear attention produce identical classifications and miss all tested long-path positives. Screening is therefore not a prerequisite for provenance collection or exact checking (Appendices~\ref{sec:certgraph_arch} and~\ref{sec:graph_evaluation}).

\subsection{Paper Organization}
Section~\ref{sec:related_conclusion} positions the work. Section~\ref{sec:threat_model} defines the threat and graph abstraction; Section~\ref{sec:cipg_provenance} develops the evidence contract. Sections~\ref{sec:game_theoretic} and~\ref{sec:evaluation} present remediation and evaluation. Sections~\ref{sec:discussion} and~\ref{sec:conclusion} discuss validity and conclusions. The appendices give the auxiliary learning study and artifact corrections.
